% =====================================================================
%  HIKET — MAIN MANUSCRIPT (working draft; renamed from HIKET_story_outline)
%  Materials & Methods written out in detail (living section, updated as the
%  analysis proceeds); the remaining parts stay sketched here and are developed
%  in the storyline note, then migrated in. Headings to be collapsed later.
%
%  PUBLIC narrative schema = Schimel OCAR (Opening / Challenge / Action /
%  Resolution), used as the four \part divisions below.
%
%  PRIVATE scaffold (NEVER surfaced in the text — author reference only):
%    Shadow      = the observed accumulation->saturation is precisely what a
%                  steady-state start CANNOT produce (an equilibrium start is
%                  already "saturated"); this triggered transient init
%    Persona     = good R-hat / physical-looking bias masking artefacts
%    Mentor      = published priors/lineage: real but boreal-blind, then freed
%    Trickster   = sigma_input, the shape-shifter that reveals the real
%                  treasure (input/history uncertainty)
%    Unconscious = the unrecoverable pre-observation history; the humus pool
%                  that "remembers" equilibrium (deep time made physical)
%    Self        = the fair, equal-DoF, exactly-integrated, homogenized frame
%    Boon        = a method + a relocated question, not a verdict on Finland
%
%  Protagonist (held constant): the INITIALIZATION PROBLEM — how to start a
%  non-equilibrium simulation with no history — with transient init as the move.
%
%  RECURRING THREAD — Bayesian inference (tagged % [Bayes] at each touchpoint):
%  state of the art, routinely applied, only mildly novel here — but it is the
%  machinery that MAKES the solution possible: it turns the unknown initial
%  state from a fixed assumption into a CALIBRATED parameter (sigma_init) whose
%  uncertainty propagates into both the initial state and the forward run.
%  Surfaces 4x: Challenge (fixed-value trap vs. calibrate-it), Action (the
%  enabling machinery + anchor propagation), Resolution (the lens that makes
%  non-identifiability/sigma_input visible), Outlook (priors from reconstructed
%  history). Keep it a thread, not a section — it must not crowd the protagonist.
%
%  RECURRING THREAD — historical perspective (tagged % [History]): a quick
%  history of (Finnish) forest management is what PHYSICALLY justifies the
%  non-equilibrium — the violation of the equilibrium assumption. General mention
%  in the Intro (own small subsubsection) -> detailed examination in the
%  Discussion -> closes the Conclusions (refinement needs historical data).
% =====================================================================

\documentclass[11pt]{article}

\usepackage[margin=2.5cm]{geometry}
\usepackage{amsmath}
\usepackage{graphicx}
\usepackage{titlesec}
% ---------------------------------------------------------------------
%  MODULAR APPENDICES.
%  Each appendix lives in appendices/*.tex as a `subfiles` document. It can be:
%    (a) compiled STANDALONE  -> `pdflatex appendices/appendix_xxx.tex`
%        (borrows this file's preamble automatically), or
%    (b) compiled TOGETHER    -> pulled in here via \subfile{...} (see the
%        Appendices part near the end). Comment a \subfile line to drop a module.
%  This keeps each methods insert self-contained and separately shareable.
% ---------------------------------------------------------------------
\usepackage{subfiles}
% ---------------------------------------------------------------------
%  OPTIONAL STRATIFIED-CALIBRATION STRAND (rendered in red).
%  A self-contained future-work line that can be ADDED if data/time allow, or
%  DROPPED entirely with a single switch while the whole storyline stays in place
%  (just smaller). Set \stratfalse to remove every \ifstrat...\fi block below;
%  while \strattrue they render in red so the strand is visually separable.
% ---------------------------------------------------------------------
\usepackage{xcolor}
\newif\ifstrat
\strattrue        % <-- \stratfalse drops the entire stratified-calibration strand
\newcommand{\stratnote}[1]{\ifstrat{\color{red}#1}\fi}  % inline red, droppable

% Number and list every heading level down to subparagraph.
\setcounter{secnumdepth}{5}
\setcounter{tocdepth}{5}

% Make \paragraph and \subparagraph display as block headings (outline view).
\titleformat{\paragraph}[block]
  {\normalfont\normalsize\bfseries}{\theparagraph}{1em}{}
\titleformat{\subparagraph}[block]
  {\normalfont\normalsize\bfseries\itshape}{\thesubparagraph}{1em}{}
\titlespacing*{\paragraph}{0pt}{\baselineskip}{0.3\baselineskip}
\titlespacing*{\subparagraph}{0pt}{\baselineskip}{0.3\baselineskip}

\title{[PLACEHOLDER TITLE]\\
  \large Starting from the past: transient initialization and the
  input-history problem in soil carbon models}
\author{Lorenzo Menichetti \and [co-authors TBD]}
\date{Working draft --- \today}

\begin{document}
\maketitle


% ---------------------------------------------------------------------------
% STATUS BLOCK -- rewritten 2026-09-02. The 2026-08-10 version listed four
% staleness items (multiplicative-normal likelihood, the wrong-in-sign
% over-prediction section, sigma_obs = 0.442 used as the total error, and MRT
% computed through transient_init). ALL FOUR ARE NOW RESOLVED in the text.
% Keep this block updated or delete it; do not let it drift again.
% ---------------------------------------------------------------------------
\begin{center}
\fbox{\begin{minipage}{0.93\textwidth}
\small
\textbf{DRAFT STATUS --- 2026-09-02.} All numbers are keyed to the calibration
completed 2026-09-01, RUN\_IDs \texttt{20260831\_1624*} (Roihu jobs
972096--972101), with the log-normal correlated-error likelihood
($\sigma_{\mathrm{tot}}=0.800$) and the $\sigma_{\mathrm{input}}$ prior centred
on the Lehtonen \& Heikkinen understorey anchor ($1.08$). All six models
converged (R-hat $\le 1.008$, ESS $\ge 1747$).
\textbf{Open:}
\textbf{(i)} \emph{the input anchor is not settled} --- the two external
estimates of the understorey share of boreal litter disagree ($27\%$ vs $8\%$),
implying $\sigma_{\mathrm{input}}$ prior centres of $1.27$ and $1.08$. A
companion calibration on the higher anchor exists and is held aside; the choice
is with A.\ Lehtonen, who authored both sources. Numbers keyed to that arm are
deliberately absent here.
\textbf{(ii)} figures are one calibration behind the text and must be rebuilt
before submission;
\textbf{(iii)} no citations or figure placements yet --- there is no
bibliography file;
\textbf{(iv)} the projections contain \emph{no} climate change (the last
twenty observed years are recycled), so every forecast statement is about
disequilibrium relaxation under stationary climate and must say so.
\end{minipage}}
\end{center}

\begin{abstract}
\noindent[PLACEHOLDER ABSTRACT.] One paragraph: steady-state initialization
is a near-universal but rarely examined convention; an equilibrium start is
already at its saturated end-state, so it cannot reproduce an accumulating
sink that is only now beginning to saturate --- the very dynamic observed in
Finnish forest soils. We test a transient alternative across six models under
one fair frame; it captures the accumulation-then-saturation as a proof of
concept, but relocates the binding uncertainty onto litter inputs ---
especially their unknown history --- pointing to historical, interdisciplinary
reconstruction as the real frontier.
\end{abstract}

\tableofcontents

% =====================================================================
\part{Opening}
% Schimel O: broad context + stakes. Lead with the GENERAL convention;
% Finland arrives as the vivid instance. (Generalization-first, confirmed.)
% =====================================================================

\section{Introduction}

\subsection{The convention: steady state at $t_0$}
A near-universal default in SOC modelling. State it plainly as common
practice, not yet as a problem.

\subsubsection{What the assumption actually claims}
That the system's past equals its present equilibrium (inputs $\approx$
outputs at the start). Make the hidden content of the convention explicit.

\paragraph{Why it is attractive}
Convenient, analytically clean, needs no historical forcing. The reasons it
became default.

\paragraph{Why it is rarely examined}
It is invisible precisely because it is standard. Set up the reader to
recognize their own practice.

% -------------------------------------------------------------------
% [TEMP notes from the new keystone sources (2026-07-16) — coarse statements +
%  refs to fold into prose later. Full PDFs + anchors in literature/README.md.]
\paragraph{[TEMP] The convention in the Finnish inventory (sources to fold in).}
\begin{itemize}
\item The steady-state start is the \emph{operational} convention of the Finnish
  national soil-C inventory: Yasso07 is run ``into a steady state'', assuming
  ``the average level of inputs and climate has remained steady over centuries''
  (Lehtonen et al.\ 2016, Geosci.\ Model Dev.\ 9:4169). $\Rightarrow$ precisely the
  assumption HIKET relaxes.
\item It rests on the \emph{absence of repeated stocks}: repeated soil sampling is
  ``the most straightforward way'' to see change but ``too expensive'' for forest
  soils, so dynamics were inferred indirectly (Palosuo 2008, Diss.\ Forestales 61).
\item It was flagged early: Peltoniemi et al.\ (2004, Glob.\ Change Biol.\ 10:2078)
  modelled soil-C change over stand age from a \emph{chronosequence} and noted soils
  may sit \emph{below} equilibrium (historical slash-and-burn), and that ``the
  long-term trend \dots\ cannot be distinguished from the measured data unless the
  measurements cover at least two rotations.'' HIKET is the successor this called
  for: repeated campaigns (VMI8/Biosoil/Komeetta) + a calibrated initial state.
  (Aleksi Lehtonen, HIKET coauthor, is acknowledged in Peltoniemi 2004 — frame the
  lineage collegially.)
\end{itemize}
% -------------------------------------------------------------------

\subsubsection{The non-equilibrium reality}
Real ecosystems are transient; they carry the imprint of disturbance,
land-use and productivity history.

\paragraph{An equilibrium start is already at the end}
A steady-state start sits at zero net change --- by construction already at
its saturated end-state. It forecloses the accumulation phase before the
simulation begins. Foreshadow the capture-failure without resolving it.

\subsubsection{Managed landscapes are rarely at equilibrium}
% [History] — general statement here; Finnish forest history is the worked
% example, examined in detail in the Discussion.
Land-use and management histories keep soils in transient states; equilibrium
is the exception, not the rule. Flag Finnish forest history as the concrete
case that justifies the violated assumption --- stated generally here, examined
in detail in the Discussion.

\subsection{Why this matters}

\subsubsection{Scientific stakes}
Initialization choices propagate into every downstream projection and into
model intercomparison.

\subsubsection{Policy stakes: the motivating case}
Finland reports its forest soil as a sink; climate policy leans on it.
Introduce Finland as the instance, not the subject.

\paragraph{The observed dynamic}
Accumulation (VMI8$\to$Biosoil), then recent hints of saturation (Komeetta,
under review). A sink caught in the act of saturating. State the phenomenon
the model must reproduce.

\subsection{Scope and contribution (preview)}
One short paragraph previewing the move: a transient initialization, tested
fairly across a model ladder. Keep it a promise, not the argument.
% CONTRIBUTION DELTA --- state explicitly (the ``relatively new, applied'' register;
% reviewers respect it if precise, and ``spin-up is old'' is the reflex to pre-empt).
% NOT transient init per se, but:
%   (i)  the unknown historical start made a CALIBRATED, uncertainty-PROPAGATING
%        parameter (sigma_init) when NO historical forcing data exist --- not a fixed
%        spin-up nor a prescribed land-use history;
%   (ii) applied HOMOGENEOUSLY across a structural ensemble as a fair-frame device
%        (lets divergence be attributed to structure vs. initialisation artefact);
%   (iii) at national-NFI scale (~350-400 plots, per-plot).
% The insight it buys: inputs/forcing history bind; structure does not.

% =====================================================================
\part{Challenge}
% Schimel C: the knowledge gap and the question. The Shadow lives here.
% =====================================================================

\section{The gap and the question}

\subsection{The saturation an equilibrium start cannot show}
The triggering insight: to reproduce a sink that is saturating, the model must
first reproduce the accumulation it is the tail of --- and an equilibrium
start, already saturated, cannot. This is what motivated transient init.

\subsubsection{Why this generalizes beyond Finland}
Any non-equilibrium system started at steady state forecloses its own
transient. Lift the problem off the case study.

\subsubsection{Artefact or structure?}
Do divergent saturation projections reflect genuine model structure, or the
initialization and calibration convention? The question the fair frame will
answer.

\subsection{The methodological gap}
No established, calibratable way to initialize a transient run when
historical forcing data are absent.

\subsubsection{The fixed-value trap}
% [Bayes]
Treating the unknown initial state as a single fixed number (equilibrium, or a
guessed value) hides a strong, untested assumption inside the result.

\subsubsection{Calibrate it, do not assume it}
% [Bayes] — first appearance of the thread; the gap is what Bayes will close
The alternative: make the historical start a parameter with a prior and let
the data and its uncertainty speak. Plant the thread; pay it off in Methods.

\subsection{Objectives and hypotheses}
\begin{itemize}
  \item[H1] Transient initialization lets the models reproduce the
        accumulation-then-saturation a steady-state start cannot. [placeholder]
  \item[H2] Residual model differences under a fair frame are structural and
        modest. [placeholder]
  \item[H3] The binding uncertainty migrates to inputs, not kinetics.
        [placeholder]
\end{itemize}

% =====================================================================
\part{Action}
% Schimel A: methods + results. The Self (fair, whole frame) is built here.
% =====================================================================

\section{Materials and methods}
% Detailed derivations are offloaded to the modular appendices (see the
% Appendices part); the M&M body states each choice and points to its appendix.

\subsection{Study system and data}

\subsubsection{Plots and SOC campaigns}
The calibration set is $512$ permanent forest plots drawn from the Finnish
National Forest Inventory (NFI) and its associated soil surveys
(BioSoil/MUSTIKKA/Komeetta lineage), spanning the boreal south--north gradient.
Each plot carries soil organic carbon stocks from up to three
temporally separated campaigns, which serve as the calibration and validation
targets: VMI8 ($1985$--$86$), BioSoil ($2006$), and Komeetta ($\sim\!2024$).
The repeated stocks --- two to three observations per plot separated by decades
--- are what make the accumulation \emph{dynamic} observable rather than a
single snapshot; this is the empirical basis for testing a transient against a
steady-state start. Plots with zero modelled litter, without an organic
F/H~layer, on peatland vegetation types, or with an implied mean residence time
$>100$~yr are excluded, as are four plots whose whole-profile stock exceeds
$250$~Mg\,ha$^{-1}$ in any campaign (see below) and eight whose modelled litter is a
single value repeated unchanged across all $39$ years of the series --- a fixed output of
the litter model where the underlying biomass record was static, rather than a measured
time series, and therefore uninformative about the temporal change under study. A fixed $20\%$ of the
calibration-ready plots is held out for validation. The
end use is the Finnish national greenhouse-gas inventory, for which Yasso is the
operational model.

\subsubsection{SOC stocks: a single homogenized basis}\label{sec:socdata}
All three campaigns are drawn from one LUKE source (a workbook maintained by
H.~Ilvesniemi) rather than from campaign-specific files processed along separate
paths, so that the differences between $1985$, $2006$ and $2024$ carry dynamics
rather than protocol. The stocks are LUKE's own, carrying both a bulk-density and
a coarse-fragment (stoniness) correction; applying LUKE's cleaning rules to this
source reproduces the official national figures exactly ($59.1$ and
$61.0$~Mg\,ha$^{-1}$ for organic\,$+$\,$0$--$40$~cm, region-weighted, in $2006$
and $2024$ over $446$ plots).

This matters because it corrects a bias we had carried. Our previous SOC
assembly applied no coarse-fragment correction to the mineral layers, and since
stoniness is volumetric it inflates mineral carbon multiplicatively while leaving
the organic layer alone --- precisely the signature we found, a near-constant
factor of $1.60$--$1.66$ across all three mineral layers in both years against an
organic-layer ratio of $0.86$. With a median coarse-fragment content of $43\%$,
$1/(1-0.43)=1.76$ accounts for the discrepancy. Campaign-specific processing had
also drifted: the $2024$ ingest was coded against the $1985$ layer protocol.

Because the models carry bulk SOC, the calibration target is a whole-profile
stock: measured organic layer, plus measured mineral $0$--$40$~cm, plus a
modelled deep tail. The tail follows an exponential depth model
$C(z)=C_0\exp(-\lambda z)$ with one $\lambda$ per soil class (profiles pooled,
$C_0$ concentrated out analytically) and $C_0$ then re-solved per profile against
its own measured layers. Crucially the tail is integrated only to
$z_{\mathrm{cap}}=\min(100\,\mathrm{cm},\,\text{depth augering reached})$, so the
target is a per-plot \emph{variable} depth matching how each profile was actually
measured; under a fixed $1$~m rule, $38$ plots that hit augering refusal at
$10$--$40$~cm were each assigned some $14$~tC\,ha$^{-1}$ of carbon below bedrock.
The extrapolation is checked against ground truth rather than assumed: the
BioSoil $40$--$80$~cm layer, measured on $501$ plots and held out of the fit, is
predicted with a median ratio of $0.96$ and a bias of $-2.5$~Mg\,ha$^{-1}$ ---
mildly conservative.

The resulting target rises gently and is internally consistent, with campaign
median stocks of $59.4$, $66.7$ and $67.4$~tC\,ha$^{-1}$. The earlier assembly
gave $63$, $102$ and $105$: a $62\%$ jump between $1985$ and $2006$ that was an
artefact of cross-campaign inconsistency rather than a dynamic, and a $2006$
value contradicting the official same-year figure by a factor of $1.7$.

\subsubsection{Litter inputs}
Plot-level litter is taken from the Tupek et al.\ dataset
(Zenodo \texttt{10.5281/zenodo.19736499}), already expressed in
tC\,ha$^{-1}$\,yr$^{-1}$ and fractionated into the AWEN quality classes required
by Yasso: a non-woody component (foliage and fine roots) plus fine- and
coarse-woody components (branches, coarse roots, stem bark, stumps). It is a
\emph{tree}-litter product covering the full living-tree biomass-component
turnover --- so belowground litter (fine \emph{and} coarse roots) is included ---
but understorey/ground-vegetation litter and aboveground tree mortality (deadwood)
are not represented, which becomes relevant to $\sigma_{\mathrm{input}}$ below. The population mean rises from $\sim\!2.1$ to a mid-2000s peak of
$\sim\!3.0$ and then plateaus near $\sim\!2.8$~tC\,ha$^{-1}$\,yr$^{-1}$, with a
strong south--north gradient across plots. This rise-then-plateau history is
load-bearing twice over: it is the forcing that drives the observed
accumulation-then-saturation (Results), and its level is what the input
multiplier $\sigma_{\mathrm{input}}$ later interrogates (Discussion).

The first year of the series, $1985$, is \emph{reconstructed rather than observed}. In the
source product the $1985$ values collapse to a median $0.060$~tC\,ha$^{-1}$\,yr$^{-1}$
against $1.398$ in $1986$ --- a factor of $23$ --- with an identical number of records in
both years, uniformly across months, plots and all three size classes. No ecological process
produces a single year at $4\%$ of the next and then recovers; the first year of the series
simply does not carry a turnover signal, whatever its origin (a plausible route is estimation
from between-inventory biomass increments, where the first year has no predecessor to
difference against, but the treatment below does not depend on that). This matters because $1985$ is $t_0$: it is
simultaneously the year the transient pre-run terminates, the year of the first SOC
campaign, and part of every litter anchor computed from the head of the series. Left
uncorrected it placed the end of the spin-up some $21\%$ below the true $1985$ litter
level. We therefore fit a linear trend to $1986$--$1990$ for each plot and litter
component and backcast one year. A plain multi-year average was rejected because litter is
rising across that window, so averaging would place $1985$ \emph{above} $1986$ and
contradict the growing-stock history used elsewhere in the initialization; the backcast
places it just below ($1.754$ against $1.838$~tC\,ha$^{-1}$\,yr$^{-1}$). The three
defensible reconstructions differ by $\sim\!3\%$ in their effect, against the $\sim\!21\%$
error being removed, so nothing downstream depends on which is chosen.

\subsubsection{Climate forcing}
Climate is gridded daily weather over $1961$--$2025$, aggregated per plot into
each model's climate modifier $\xi$. Because $\xi$ is a nonlinear function of
temperature (and, in Yasso15/20, precipitation), the annual modifier is
evaluated year by year and then averaged --- $\xi(\overline{\text{climate}})
\neq \overline{\xi_{\text{annual}}}$, a Jensen-inequality bias of $5$--$15\%$ in
boreal settings that we correct explicitly rather than aggregate climate first.

\subsection{The model ladder}
We span six models on two crossed axes --- pool complexity and climate
integration --- so that no conclusion rests on a single structure. The roles are
asymmetric by design and stated up front: \textbf{Yasso is the primary,
operational model} (it is the Finnish GHG-inventory model, so it is the one that
matters for the application), while \textbf{SP1/TP2/TP3 form a structural
robustness ladder}. The simpler models are not competitors for a ``best'' model
but a sensitivity test --- a check that the transient-initialization and
input-uncertainty conclusions survive across model structure. This reframes the
complexity axis as robustness rather than an optimal-complexity search.
% MODEL ROLES (state up front): Yasso is the PRIMARY / operational model --- it is
% the Finnish GHG-inventory model, so it is the one that matters for the application.
% SP1/TP2/TP3 are a structural ROBUSTNESS ladder: hypothesis-testing that the
% transient-init + input-uncertainty conclusions survive across model structure,
% NOT a beauty contest for the ``best'' model. This reframes the complexity axis as
% robustness/sensitivity rather than an optimal-complexity search.

\subsubsection{Pool-complexity axis}
Four structures of increasing pool count, all linear donor-controlled
compartment systems: \textbf{SP1}, a single bulk SOC pool (one time-scale);
\textbf{TP2}, an active and a humus pool in the ICBM two-pool form (fast litter
turnover feeding a slow reservoir through a humification fraction);
\textbf{TP3}, a three-pool active$\to$slow$\to$humus sequential cascade; and
\textbf{Yasso07}, the five-pool AWEN+humus structure in which litter is split by
chemical quality. Climbing the ladder adds internal structure while the
forcing, data, and likelihood are held fixed.

\subsubsection{Climate-integration axis}
Three five-pool Yasso variants that share the AWEN+humus structure but differ in
how climate enters the decomposition rates: \textbf{Yasso07} (temperature
response of Tuomi et al.\ 2009), \textbf{Yasso15} (recalibrated, with an
explicit precipitation term), and \textbf{Yasso20} (the Viskari et al.\ 2022
structure). For Yasso20 we free all twelve inter-pool transfer fractions
(previously six fixed as structural zeros and three algebraically derived); the
original constraints reflected a warm-site global dataset, and fixing them here
would confound climate-integration differences with a calibration asymmetry.

\subsection{Bayesian calibration: the enabling machinery}
% [Bayes] — the thread's payoff. State of the art, routinely applied; here it
% is what makes transient init possible. Implementation detail, but fundamental.
Each model is calibrated in a Bayesian framework (the \texttt{BayesianTools}
DEzs differential-evolution sampler): five chains of $50{,}000$ iterations each,
the first $5{,}000$ discarded as burn-in, with a fixed seed for
reproducibility. The posterior combines the log-normal, correlated-error likelihood
(below) with the homogenized priors; inference returns full joint posteriors
over rates, fractions, and the two auxiliary uncertainty parameters. The role of
the Bayesian frame here is enabling rather than novel: it is the machinery that
turns the unknown initial state from a fixed assumption into a calibrated
quantity.

\subsubsection{Why a Bayesian frame is the natural fit}
The initialization problem is fundamentally one of \emph{incomplete
information} about the past: we have no record of the forcing that produced the
pre-observation state. A Bayesian frame is the natural fit because the prior
encodes what little is known about that history and the posterior carries the
residual uncertainty forward, instead of collapsing it to a single assumed
value as the steady-state convention does.

\subsubsection{Uncertainty that propagates, not a likelihood patch}
$\sigma_{\mathrm{init}}$ is a prior on the strength of the historical (initial)
carbon state, not a term added to the likelihood at the end. Its spread
therefore propagates into \emph{both} the initial condition and the forward
dynamics: a wider $\sigma_{\mathrm{init}}$ legitimately widens the
posterior-predictive envelope because it represents genuine ignorance about the
starting point, not observation noise. This distinction --- a propagated prior
on the state versus a fitted error term --- is what makes the transient start an
inferential object rather than a numerical fix.

\subsection{Transient initialization (the protagonist)}

\subsubsection{Pushing the equilibrium assumption back in time}
Rather than assume equilibrium at the observation start $t_0$ (VMI8, $1985$),
each draw initializes the pools at their analytical steady state under a much
earlier historical flux (the pre-initialization year $1917$) and then runs the
model \emph{transiently} forward from $1917$ to $t_0$. The equilibrium
assumption is thus displaced $\sim\!70$~years into the past, where it is far
less consequential: by $t_0$ the imprint of that assumed start has largely
dissipated, and the model arrives at the observation window still accumulating
--- the state an equilibrium start at $t_0$ can never occupy.

\paragraph{The fast/slow relaxation argument}
The pre-run works because the pools relax on separated time-scales. The fast
pools turn over in years and fully forget their initial value within the
$68$-year spin-up; only the slow humus pool, with a residence time of a century
or more, still ``remembers'' the assumed start at $t_0$. Displacing the
equilibrium assumption backwards therefore confines its residual influence to
the one pool where it is least resolvable --- and it is exactly this pool whose
uncertain strength $\sigma_{\mathrm{init}}$ absorbs.

\subsubsection{Calibrating the historical anchor}
% [Bayes]
The pre-observation forcing is not known, so it is not fixed. The strength of
the $1917$ historical state enters as the calibrated parameter
$\sigma_{\mathrm{init}}$ (a proportional scaling on the initial carbon state,
log-transformed, prior centred at $1$), which every draw infers jointly with
the rest. The \emph{shape} of the input path between $1917$ and $t_0$ is no
longer a bare linear ramp: it follows the reconstructed NFI growing-stock
history, since
litter scales with standing biomass --- a physically grounded, near-free
reconstruction of the pre-observation input trajectory that targets the coarse
pre-$1985$ phase (see the forest-history discussion and the growing-stock
figure). Full pre-run construction in Appendix~\ref{app:transient}.

\subsection{A fair frame for intercomparison}
The design principle: any residual difference must be attributable to
structure alone.

\subsubsection{Shared data, likelihood, forcing, filtering}
All six models see exactly the same inputs: the same $512$ plots, the same
litter and climate forcing, the same likelihood and error structure, the same
plot-exclusion filtering, the same holdout split, and the same fixed seed. Any
residual difference between models is therefore attributable to structure, not
to a difference in what each model was shown.

\subsubsection{Exact integration in every model}
All six models integrate their linear pool system \emph{exactly}, so that
differences between them are structural rather than numerical: SP1 and TP2 use
the closed-form per-step solution; TP3 uses the closed-form matrix exponential
of its lower-triangular cascade generator (divided differences of the pool
eigenvalues); and the Yasso variants use a Fortran matrix-exponential
back-end. This matters because TP3 was previously the sole approximate
integrator --- an explicit forward-Euler step that rings when the active pool
turns over faster than the annual step ($\alpha_A\xi>1$), producing an
interannual oscillation once mistaken for a climate response. With exact
integration that numerical artefact is removed; a bounded, genuinely physical
fast-pool oscillation survives (Discussion). Solvers and the discarded Euler
path in Appendix~\ref{app:exact}.

\subsubsection{Homogenized priors: method, not number}
Priors are homogenized in \emph{construction, scale, constraint, and degrees of
freedom} --- not in raw value (identical numbers, tried first, broke Yasso07).
A three-tier scheme applies to every model: (i)~climate and size parameters take
per-model centres and widths from that model's \emph{own} published calibration,
homogeneous in method and scale rather than magnitude; (ii)~every inter-pool
partitioning fraction --- the twelve lateral transfer fractions in the Yasso
family and the one or two sequential humification fractions in TP2/TP3 --- shares
a common weakly-informative logit prior (SD~$0.4$) with per-model external
centres, wired by explicit per-fraction listing; and (iii)~the auxiliary
parameters $\sigma_{\mathrm{init}}$ and
$\sigma_{\mathrm{input}}$ use an identical construction everywhere. Because the
construction is common, any surviving divergence is attributable to structure
and data, not to prior asymmetry. Full parameter-class $\times$ prior-criteria
tables in Appendix~\ref{app:priors}.

\subsubsection{Degrees of freedom, set by external information}
% ROUND-1 REVISION (finalized + validated 2026-07-16): the simple-model kinetics
% are homogenized with HOW YASSO IS ALREADY PARAMETERIZED. Yasso does not fit its
% decomposition rates --- the rate vector is held fixed at its litterbag values;
% what Yasso calibrates is its inter-pool (partitioning) fractions, climate, size,
% and the two sigmas. The simple models now mirror that: intrinsic rates externally
% anchored (fast fixed, slow via a very-informative prior), partitioning
% fractions free. Anchor validated in doublechecks/icbm_anchor_sanity.R (sigma_input
% lands 1.05-1.36, inside the physical understorey window). See NEXT\_SESSION.md §5 (C1).
The fair frame requires not just equal degrees of freedom but equal external
\emph{information}, and the cleanest way to secure it is to parameterize every
model the same way Yasso already is: \emph{its decomposition rates are not fit to
the soil-carbon data at all}. In Yasso the rate vector is held fixed at its
litterbag-calibrated values (Tuomi et al.\ 2009/2011); what Yasso calibrates is
the twelve inter-pool transfer fractions, the climate response, the woody-size
modifier, and the two auxiliary uncertainty parameters. The simple models are
brought onto the same footing. Their \emph{intrinsic decomposition rates} are
anchored to the Ultuna bare-fallow ICBM calibration (Andr\'en \& K\"atterer 1997;
$k_1=0.8$, $k_2=0.00605$~yr$^{-1}$): the fast rate --- litterbag-grounded and
well transferable --- is held fixed, while the single poorly-transferable slow
rate ($k_2$, optimized on arable soil) carries a very-informative prior, so the
data may nudge it and the posterior width reports how far a boreal forest soil
departs from that anchor. Their \emph{humification fractions}, by contrast, are
the analogue of Yasso's free lateral fractions --- sequential partitioning of
outflow into the next, slower pool --- and are therefore left \emph{free} under
the same common logit prior, re-centred on the ICBM humification value. Fixing
them too would leave the simple models with less partitioning freedom than Yasso,
not more homogeneity. How completely a structure's rates can be anchored is set by
whether its time-scale structure matches an existing calibration: TP2 (two
time-scales) maps onto ICBM one-to-one; SP1 (one time-scale) inherits the ICBM
bulk residence time; and TP3 (three time-scales) places both of its slow pools at
the ICBM ``old'' subsystem scale, its fast pool at the young scale. The result is
uniform across all six models: \emph{no model fits its kinetics to the SOC data},
and every model spends its calibrated freedom on the same four things --- the
partitioning fractions, the climate response, the litter input, and the initial
state. Climbing the pool-complexity ladder therefore adds only free partitioning
degrees of freedom (none in SP1, one in TP2, two in TP3, twelve in the Yasso
family); identifiability then tracks information, by design rather than by
accident.
\stratnote{[strand] This convention is also what makes the stratified extension
Yasso-only: with rates anchored, Yasso's stratifiable knobs are inputs + the
humus-formation fraction --- see the optional strand in the Outlook.}

\subsection{Calibration and inference}

\subsubsection{Error model}
% ⚠ REWRITTEN 2026-09-02. This described a multiplicative-normal likelihood and
% promised a Discussion section attributing an ~11% over-prediction to it. BOTH
% are superseded: log-normal has been the default since 2026-08-07, the correlated
% error structure since 2026-08-19 (SLURM: HIKET_CORRELATED_LIK=1,
% HIKET_SIGMA_TOT=0.800, HIKET_SIGMA_1985_INFL=1), and the over-prediction is gone.
Residuals are modelled on the log scale, where they are near-Gaussian and where
the error is proportional to the stock --- the appropriate form for SOC, which
spans an order of magnitude across plots.

The scale of that error is not the measurement error. The homogenisation of the
SOC campaigns yields a measurement CV of $0.43$, but the observed spread of log
residuals is close to $0.8$ in every model; the difference is model error, and
it is the larger of the two. We therefore set the \emph{total} error scale to
$\sigma_{\mathrm{tot}} = 0.800$ rather than using the measurement CV in its
place, and verify afterwards that the fitted residual spread reproduces it.

Crucially, that total is \emph{partitioned} rather than treated as independent
noise. Plot-years are not independent observations: repeated measurements of the
same plot correlate strongly across campaigns, plots within a latitude band share
a level, and each campaign carries its own comparability offset. Treating $1205$
plot-years as independent overstates the information they carry about the
national \emph{level} --- which is precisely the quantity that pins
$\mathrm{MRT}\times\sigma_{\mathrm{input}}$. We therefore write
\[
  \log y_{ij} \;=\; \log f_{ij}(\theta) \;+\; u^{R}_{r(i)} \;+\; u^{P}_{i} \;+\; u^{C}_{j} \;+\; e_{ij},
\]
with a regional offset $u^R$, a plot offset $u^P$, a campaign offset $u^C$ and
independent noise $e$. No variance component is estimated: all four are fixed
from the observations with campaign levels removed --- $\tau_R = 0.117$ from
latitude-band means, $\tau_P = 0.396$ from the $2006$--$2024$ pair covariance,
$\tau_C = 0.06$ for $1985$ and $0.03$ for $2006$ and $2024$, and
$\sigma_e = 0.685$ as the remainder of the unchanged total. The offsets are
marginalised analytically, so no parameters are added; because every variance is
fixed, the covariance matrix is constant and is factorised once. Setting all
$\tau$ to zero reproduces the independent log-normal likelihood exactly, which is
the unit test the implementation is held to.

Two consequences are worth stating plainly. First, this is a decision about how
much the data constrain the level, and $\tau_C = 0$ --- the implicit choice in a
conventional likelihood --- is also a decision, and the least defensible one.
Second, the correction removes information about the \emph{level} while leaving
within-plot contrasts untouched, so it loosens what the level pins and should not
be expected to move quantities fixed by trajectory shape. Log-likelihoods are not
comparable across this change, since the normalising constant moves; models are
compared on predictive skill, not on likelihood.

\subsubsection{Physically-bounded input priors}
% Now a METHODS choice (implemented for the production re-calibration), not only
% an Outlook refinement. The effective litter flux is bounded to the boreal
% NPP envelope, homogeneous across models, via a bounded re-parameterisation.
Left as a scale-free multiplier centred at one, $\sigma_{\mathrm{input}}$ can
drag the effective litter flux to physically impossible levels (the unbounded
prior let two models reach $13$--$20\times$ boreal net primary production). We
therefore bound the \emph{effective flux} --- $\sigma_{\mathrm{input}}\times J$,
and the coupled $1917$ pre-run flux --- to a boreal NPP envelope
($\sim\!0.5$--$8.7$~tC\,ha$^{-1}$\,yr$^{-1}$), homogeneous across models, via a
bounded re-parameterisation of the $\sigma_{\mathrm{init}}/\sigma_{\mathrm{input}}$
pair (a scaled-logit ``flux-pair'' transform) rather than a hard truncation, so
the sampler still mixes across the boundary. In addition, because the litter
product, though it includes belowground (root) litter, structurally omits
understorey/ground-vegetation and mycorrhizal inputs (and aboveground tree
mortality), the $\sigma_{\mathrm{input}}$ prior is centred on the physically
expected correction ($>1$) rather than on one, so its posterior is read against
expected total litter rather than against the tree-only baseline.
Rationale, envelope and mechanism in Appendix~\ref{app:input-bounds}.

\subsubsection{Sampler and diagnostics}
Sampling uses the DEzs differential-evolution sampler (five chains,
$50{,}000$~iterations, $5{,}000$~burn-in). Mixing is checked by the
Gelman--Rubin $\hat{R}$ (target $<1.05$) and effective sample size, and the
proposal infinity-rate is monitored to catch prior-driven forward blow-ups. We
distinguish mixing from \emph{identifiability}: with healthy mixing the sampler
traverses the non-identified ridges cleanly, so a fraction can show good
$\hat{R}$ yet remain data-uninformed. The genuine non-identifiability signature
is therefore read elsewhere --- prior-pinning, near-perfect same-pool
anti-correlation ridges, and low KL divergence of posterior from prior
($<1$~nat) --- not from poor $\hat{R}$.

\subsubsection{Holdout split and validation}
A fixed subset of plots is held out of the likelihood entirely and used only for
out-of-sample evaluation. Predictive skill (R$^2$, RMSE, bias) is reported
separately for the calibration and holdout sets, so that in-sample fit and
genuine predictive performance are never conflated --- the distinction that
later separates in-sample skill from the forecast divergence.

\subsubsection{Software, compute and reproducibility}
The pipeline is implemented in R, with the Yasso decomposition solved by a
compiled Fortran matrix-exponential back-end and the simple models solved in
closed form. All stochastic steps use a fixed seed (\texttt{set.seed(2025)}),
and the production calibration runs on an HPC cluster (one job per model). Code
is version-controlled and the final outputs are archived openly; per-model
posterior and predictive bundles are keyed by run identifier so that each
predictive run is traceable to its calibration.

\subsection{Predictive and residual analysis}
For each model we draw posterior-predictive SOC trajectories (an ensemble of
parameter draws pushed through the forward model, initial state included), from
which the campaign-year stocks, skill metrics, and forecast envelopes are taken.
Plot-level residuals are then regressed against candidate site covariates
(soil chemistry, texture, climate summaries, stand basal area, litter quality)
by random forest, to locate where recoverable signal remains after calibration.
Across models the leading, shared driver is stand basal area --- a productivity
and input proxy, not a kinetic variable --- pointing the refinement toward
better inputs rather than more pool structure.

\section{Results}
% =====================================================================
% WRITTEN 2026-09-02, replacing the bullet skeleton. Every number here is
% keyed to RUN_IDs 20260831_1624* (Roihu jobs 972096-972101) and was pulled
% from the posteriors and predictive bundles directly, not from the figure
% PNGs, which are still one calibration behind. Provenance per subsection is
% given in the comment above it. When the figures are rebuilt, CHECK THESE
% NUMBERS AGAINST THEM rather than assuming they agree.
%
% Observed rates and levels use the DECIDED reporting basis (Lorenzo,
% 2026-08-17): balanced plot set (observed in all three campaigns, n = 310),
% unweighted, whole profile, TRUE observation years. build_F5_stock_change.R
% was brought onto that basis on 2026-09-02; before then it used its own
% campaign cutoff (which filed the 65 plots dated 1995 as 2006 observations)
% and a nominal 39-yr denominator for a span that is really ~35 yr.
% =====================================================================

\subsection{All six models converged}
% Provenance: <MODEL>_rhat_20260831_*.txt, <MODEL>_ess_20260831_*.txt.
Every model converged cleanly. Across the six, the largest univariate
Gelman--Rubin statistic is $1.008$ and the largest multivariate potential scale
reduction factor is $1.028$ (Yasso20); the smallest effective sample size for
any parameter is $1747$, from five chains of $50{,}000$ iterations with the
first $5{,}000$ discarded. No chain contained a non-finite log-likelihood.

This matters more than a routine diagnostic statement, because the natural
expectation for these models is the opposite. With of order twenty free
parameters and two SOC observations per plot, the transfer fractions are
structurally under-determined, and it is tempting to read that as something
that should show up as poor mixing. It does not, and we return to why in
Section~\ref{sec:identifiability}: clean convergence is what a well-mixed
sampler does on a non-identified ridge, and it is evidence about the sampler,
not about the information in the data.

\subsection{Capturing accumulation-then-saturation}
% Provenance: modelled vs observed campaign means, balanced n=310, computed
% from <MODEL>_residuals_20260831_*.csv aggregated by campaign_of(year).
% Observed campaign levels from obs_campaigns(): 64.5 / 71.5 / 73.6.
The observations describe accumulation that is beginning to level off. On the
balanced plot set the mean stock rises from $64.5$~tC\,ha$^{-1}$ at the first
campaign (mean year $1989$) to $71.5$ in $2006$ and $73.6$ in $2024$: a gain of
$7.0$~tC\,ha$^{-1}$ over the first interval and $2.1$ over the second. That is
the shape an equilibrium start cannot produce, because an equilibrium start is
already at the end of it.

Started transiently, all six models reproduce the accumulation
(Figure~[F3]). None sits more than $3.5$~tC\,ha$^{-1}$ from the observed level
at any campaign, and every one of them rises across the first interval. The
\emph{shape} of the later trajectory, however, separates them into three
groups, and the separation is structural rather than a matter of skill:

\begin{itemize}
\item \textbf{Accumulation then flattening} --- Yasso07 ($66.1 \to 74.9 \to
  75.0$) and Yasso15 ($66.9 \to 76.4 \to 75.5$) reproduce the bend-over,
  rising steeply to $2006$ and then approximately holding.
\item \textbf{Accumulation without saturation} --- TP2 ($64.2 \to 73.0 \to
  76.3$) and TP3 ($64.3 \to 72.9 \to 76.1$) match the early levels most
  closely of any model but keep climbing, so they end above the observations
  having passed through them.
\item \textbf{Accumulation then decline} --- SP1 ($67.4 \to 75.7 \to 71.5$) and
  Yasso20 ($68.0 \to 74.7 \to 71.3$) overshoot at $2006$ and then turn down,
  ending below the observed $2024$ level.
\end{itemize}

\subsubsection{The rate of change, which is what an inventory reports}
% Provenance: build_F5_stock_change.R, rebuilt 2026-09-02 on the balanced set
% with per-plot true intervals. Observed +0.259 +/- 0.040 (1985-2024) and
% +0.117 +/- 0.066 (2006-2024) reproduce the decided basis exactly.
% ⚠ This SUPERSEDES the previously recorded picture ("over the full span every
% model is a sink but 1.6-2.5x too strong"), which was measured on the
% contaminated campaign mapping and the nominal denominator.
Stocks are what these models predict; \emph{changes} in stock are what a
greenhouse-gas inventory reports, and they are a far more demanding target
(Figure~[F5]).

Over the full span the ensemble brackets the observations and two members land
on them. Against an observed $+0.259 \pm 0.040$~tC\,ha$^{-1}$\,yr$^{-1}$,
Yasso07 gives $+0.258$ and Yasso15 $+0.250$ --- agreement to within a hundredth
of a unit, and well inside the observational standard error. TP2 ($+0.347$) and
TP3 ($+0.339$) run about a third too strong; SP1 ($+0.124$) and Yasso20
($+0.097$) about half too weak. The spread across six structures, $+0.10$ to
$+0.35$, is roughly six times the standard error on the observed value: the
ensemble contains the answer but cannot identify it.

Over the recent interval the picture is different, and the reason is as much
observational as structural. The observed $2006$--$2024$ rate on this basis is
$+0.117 \pm 0.066$~tC\,ha$^{-1}$\,yr$^{-1}$ --- positive, but with a confidence
interval that includes zero. The models spread from $+0.182$ (TP2) to $-0.237$
(SP1), with Yasso20 also negative ($-0.191$) and Yasso07 and Yasso15 close to
flat ($+0.010$, $-0.048$). Two models therefore report a source over the recent
period. We resist the obvious reading that they are simply wrong: over an
eighteen-year interval the observations themselves cannot establish a sink at
conventional confidence, so this interval discriminates far less than its
apparent precision suggests, and the campaign-comparability corrections of
Section~\ref{sec:socdata} bear on it much more heavily than on the full span.

\subsection{Structural differences are a weak lever in-sample}
% Provenance: predictive bundles, metrics_calib / metrics_holdout.
% ⚠ metrics_calib is known to include held-out observations in the multimodel
% builder; the values here are read from the bundles, where T1 also reads them.
Under the fair frame --- same plots, same forcing, same likelihood, same
priors by construction, same holdout --- the six models are close to
indistinguishable in fit (Table~[T1]).

Calibration $R^2$ ranges from $0.016$ (TP3) to $0.021$ (Yasso15). On the
holdout it is indistinguishable from zero in every model, the largest being
$0.0012$. Root-mean-square error on the holdout spans $46.9$ to
$48.6$~tC\,ha$^{-1}$, a range of $3.6\%$, against stocks near
$70$~tC\,ha$^{-1}$. Median bias runs $+1.0$ to $+3.3$~tC\,ha$^{-1}$ on the
calibration plots and $-0.9$ to $+0.7$ on the holdout.

Two features of that table deserve emphasis. First, \textbf{complexity buys
essentially nothing}: SP1, with one pool and six free parameters, and Yasso15,
with five pools and twenty-six, differ by $0.005$ in calibration $R^2$ and by
$0.1$~tC\,ha$^{-1}$ in holdout RMSE. The best and worst holdout RMSEs belong to
a five-pool and a three-pool model respectively; the ordering does not follow
structure. Second, \textbf{none of the models has meaningful cross-plot skill}.
An $R^2$ near zero on held-out plots means that, given a plot's litter input and
climate, these models cannot say whether its soil carbon is high or low. This is
not a failure of any one structure --- it is shared by all six, which is what
makes it informative.

\subsection{What the data can and cannot constrain}\label{sec:identifiability}

\subsubsection{Convergence without identifiability}
% Provenance: KL/marginal diagnostics per model (F7, F8, F8b, S11).
The paradox flagged above --- clean convergence in a model whose fractions
cannot be determined from two observations per plot --- resolves as follows.
The transfer fractions converge in R-hat and are simultaneously not learned from
the data: their posteriors sit within a fraction of a prior standard deviation
of their prior centres, they carry very little information relative to those
priors, and fractions leaving the same pool are almost perfectly
anti-correlated. That last signature is the mechanism --- the SOC observations
constrain the net flux out of a pool, not how it is divided --- and it is why
good mixing and weak identification coexist. A sampler that traverses a flat
ridge cleanly will report excellent diagnostics; poor R-hat would indicate a
mixing problem, not an information problem, and conflating the two would have
led us to fix the wrong thing.

\subsubsection{The input multiplier does the work}
% Provenance: posterior medians of sigma_input/sigma_init; Jbar = 2.511
% (flux_pair, confirmed in the launch logs). NB build_T1_model_summary.R still
% hard-codes Jbar = 2.472 and should be corrected to 2.511.
What the data \emph{do} constrain is the level, and the level is carried by the
input multiplier. Posterior medians of $\sigma_{\mathrm{input}}$ run $1.27$
(SP1) to $1.73$ (TP3); multiplied by the shared litter baseline $\bar J =
2.511$~tC\,ha$^{-1}$\,yr$^{-1}$ these are effective litter fluxes of $3.19$ to
$4.34$~tC\,ha$^{-1}$\,yr$^{-1}$. Every model raises its inputs above the
tree-litter baseline, and the ordering is not by complexity: the one-pool model
asks for the least and the three-pool model for the most, with the three Yasso
variants in between.

The historical anchor behaves the same way. $\sigma_{\mathrm{init}}$ --- the
ratio of $1917$ litter input to $1985$ litter input --- takes posterior medians
of $0.49$ (TP3) to $0.84$ (SP1), placing the implied $1917$ carbon stock at
$41.9$ to $52.3$~tC\,ha$^{-1}$ and the implied pre-observation accumulation rate
at $+0.05$ to $+0.19$~tC\,ha$^{-1}$\,yr$^{-1}$, in every case below the observed
modern rate. Five of the six models cluster tightly ($0.49$--$0.59$); SP1, which
has only one pool and therefore no way to hold a legacy separately from the
active store, is the outlier.

\subsubsection{Turnover, and the ridge it shares with the inputs}
% Provenance: doublechecks/intrinsic_mrt.rds, stamped 20260831_162457
% (provenance verified). ⚠ Use intrinsic_mrt.R ONLY -- the engine binding named
% steady_state is transient_init and is contaminated by sigma_init.
% ⚠ TWO published-MRT bases exist. These are the published POINT values.
For the three Yasso variants a published parameterisation exists, so the
calibrated turnover can be compared with an external estimate of the same
quantity. Intrinsic mean residence time --- a property of the generator alone,
computed at a fixed reference climate and litter composition, independent of
both auxiliary parameters --- is $24.6$~yr for Yasso07 ($90\%$ interval
$20.7$--$29.1$), $27.1$ for Yasso15 ($23.6$--$30.9$) and $22.6$ for Yasso20
($19.6$--$26.4$). The published point values are $33.5$, $30.4$ and $19.0$~yr.
Yasso15 overlaps its published posterior; Yasso20 sits entirely above its
published point; Yasso07 falls well short of its own.

These displacements should not be read as three independent findings, and in
particular should not be added to the input result above. A soil carbon stock
constrains the \emph{product} of turnover time and input rate. Faster turnover
with larger inputs and slower turnover with smaller inputs fit the observed
level equally well, and the posterior explores that trade-off directly: within
the Yasso family the two are strongly anti-correlated, so ``inputs too high''
and ``turnover too fast'' are one under-determination described from two ends.
Where a model lands along that ridge is set by the priors, and we treat the
displacement as a single result reported once.

\subsection{Residual structure: the missing predictor is not kinetic}
% Provenance: <MODEL>_rf_summary_20260831_*.rds (oob_r2) and
% <MODEL>_rf_importance_20260831_*.csv. n = 954 calibration observations.
% ⭐ This is the sharpest result in the section: predictability EXISTS in the
% data (RF explains 44-65% of the residuals) but no model captures it, and the
% covariate that carries it is an input proxy, not a kinetic variable.
The models explain about two percent of the variance in observed SOC. A random
forest trained on site covariates explains between $44\%$ and $65\%$ of what
they leave behind (out-of-bag $R^2$, $n = 954$; $52\%$--$70\%$ on the holdout
residuals). The predictability is therefore present in the data. No model
recovers it.

The covariate carrying most of it is the same in all six: stand basal area at
the first campaign, followed by development class and coarse-fragment content.
Basal area is a productivity and standing-biomass proxy --- an \emph{input}
variable --- not a kinetic one, and the fact that six structurally unrelated
models leave the same input-shaped signal in their residuals is the strongest
in-sample evidence in this study that the binding limitation is the
representation of carbon entering the soil rather than the representation of its
decay. The importance is diffuse: no single covariate dominates, which is
consistent with a broad productivity gradient the models see only through a
single national litter multiplier (Figure~[F11]).

% =====================================================================
\part{Resolution}
% Schimel R: discussion + significance + the opened frontier. The Boon.
% =====================================================================

\section{Discussion}

\subsection{Capturing saturation requires remembering the accumulation}
Resolve the Challenge: the difficulty was the equilibrium start, not the soil;
a transient start lets every model represent the observed dynamic. The central
payoff.

\subsubsection{Proof of concept, not a verdict on Finland}
What the result does and does not license. Calibrate the claim.

\subsubsection{Generalization to the model class}
Wherever steady-state init meets a non-equilibrium system, it forecloses the
transient the data demand. Lift off Finland.

\subsection{The binding uncertainty is the inputs (the $\sigma_{\mathrm{input}}$ reveal)}
% Distinct beat (chosen: clearer over woven). Trickster pays off here.
The data cannot resolve kinetics; the input multiplier absorbs what structure
cannot. So the frontier is inputs, not soil chemistry.
% CENTRALITY: not a side story --- it is the ANSWER to the complexity question
% (structure is a weak lever BECAUSE inputs dominate). Give it Discussion-central
% status, not a sidebar.
% CLOSED LOOP: diagnosed on the unbounded run -> bounded in Methods -> consequence
% on the bounded re-run (whether TP2/TP3 take a visible R^2 hit, or the misfit
% relocates onto the now-pinned kinetics). The bound is what gives the reveal teeth.
% HEDGE (honesty): say ``consistent with input/forcing error'', NOT ``proven'' ---
% the current global design cannot exclude local kinetic variation (see excursus +
% the optional stratified strand, which is exactly the fair test of that).

\subsubsection{The multiplier has a physical referent (quantitative payoff)}
% The Trickster's tell: sigma_input is NOT a nuisance -- it converts to a flux
% that external litter budgets can be laid against. Effective flux =
% sigma_input x Jbar, with Jbar = 2.511 tC/ha/yr shared by all six.
%
% ⚠ REWRITTEN 2026-09-02. The previous version argued from a bounded-vs-unbounded
% counterfactual ("the ceiling costs skill for ONLY TP2/TP3, which is the cleanest
% proof that inputs were the lever"). THAT EXPERIMENT WAS NEVER RUN, and the hard
% flux window [0.05, 4.62] was deliberately abandoned (2026-08-31): on the current
% calibration nothing is absurd, and bounding would IMPOSE the input conclusion
% rather than find it. The old text also quoted TP2 at 34 and TP3 at 50 tC/ha/yr
% from a calibration three runs back. Both are gone. What survives -- and it is the
% part that was always doing the work -- is that sigma_input is dimensionally a
% litter flux and can therefore be confronted with independent budgets.
The reason this parameter is worth reading at all is that it is not scale-free.
Multiplied by the shared litter baseline
($\bar J = 2.511$~tC\,ha$^{-1}$\,yr$^{-1}$), $\sigma_{\mathrm{input}}$ becomes an
\emph{effective litter flux to the soil} --- a quantity that independent budgets
can be laid against. Calibration places all six models in a narrow band,
$3.19$--$4.34$~tC\,ha$^{-1}$\,yr$^{-1}$ (SP1 lowest, TP3 highest), with the
complexity ordering conspicuously absent: the one-pool and five-pool models
bracket each other rather than separating.

That band sits above every independent estimate of Finnish litter input, and by
a margin that is modest rather than absurd. Against the national total litter
including understorey of Lehtonen and Heikkinen ($\sim\!2.7$~tC\,ha$^{-1}$\,yr$^{-1}$),
the six models want $18$--$61\%$ more. Against the Nordic mean aboveground NPP
compiled by Gower and colleagues ($\sim\!3.2$~tC\,ha$^{-1}$\,yr$^{-1}$), SP1,
Yasso07 and Yasso15 sit comfortably inside while TP2, TP3 and Yasso20 sit
$30$--$35\%$ above it --- tolerable only if Finnish productivity runs above that
compilation's mean. No model exceeds the Nordic \emph{maximum} ($4.62$), so the
comparison that is comfortable is the one against NPP, and the comparison that
is uncomfortable is the one against litter. We state both.

We deliberately did not close this gap by fiat. A hard ceiling on the effective
flux was specified and then not imposed, because imposing it would have written
the conclusion into the prior: the unbounded calibration is the better evidence
that these models want more input than the published budgets supply. The prior
window remains wide enough that the posterior is located by the data and the
litter-based centre, not by a wall.

Two cautions bound how far this can be pushed. First, the discrepancy is
\emph{consistent with} input or forcing error, not proof of it: a global
calibration cannot exclude local kinetic variation (see the excursus, and the
stratified strand that would be its fair test). Second --- and this is the
constraint that governs the whole Discussion --- ``inputs too low'' and ``turnover
too fast'' are one under-determination viewed from two ends, not two independent
findings. The level of a soil carbon stock constrains the \emph{product}
$\mathrm{MRT}\times\sigma_{\mathrm{input}}\times\bar J$; where that product is
split between the two factors is set by the priors, not by the SOC data. We
therefore report the displacement once, and do not bank it twice.

\subsubsection{Historical inputs: the unrecoverable past}
No forcing data pre-observation; $\sigma_{\mathrm{init}}$ stands in for an
entire history, held in the slow pool that remembers it. The deepest
uncertainty.

\subsubsection{Belowground inputs and other ramifications}
Root/rhizodeposition inputs poorly known; a further axis of the same problem.

% -------------------------------------------------------------------
% [TEMP notes from the new sources (2026-07-16) — fold into prose later.]
\paragraph{[TEMP] Missing understorey litter (why $\sigma_{\mathrm{input}}$ > 1 is expected).}
\begin{itemize}
\item The litter product used here is \emph{tree litter} --- it covers the full
  living-tree biomass-component turnover, so \emph{roots (fine and coarse) are
  included} (confirmed by B.\ Tupek and the dataset size-class table:
  \texttt{nwl}=foliage+fine roots, \texttt{fwl}=branches+coarse roots+stem bark,
  \texttt{cwl}=stumps). What is \emph{not} represented, in order of magnitude:
  understorey/ground-vegetation (incl.\ moss) litter (dominant); mycorrhizal
  mycelial turnover/exudates (large, uncertain); and aboveground tree mortality
  (deadwood), a genuine omission but small in managed forest. So a multiplier above
  one is a physical expectation, not manufactured carbon --- hence the prior is
  re-centred at $\sim$1.3 (understorey-dominated) rather than 1.
\item Grounding: Lehtonen et al.\ (2016) found understorey litter \emph{underestimated}
  in Yasso07, especially in the north (bryophytes 20--60\% of NDVI at high latitude);
  Finnish litterfall studies put understorey at $\sim$15\% (south) to $\sim$33\% (north)
  of aboveground litter. Peltoniemi et al.\ (2004) give component turnover rates for
  dwarf shrubs/bryophytes/herbs/lichen (S.\ Finland).
\item Caveat: the fraction rises strongly S$\to$N, so a single national multiplier cannot
  capture it --- an argument for stratifying inputs by site class (future work).
\end{itemize}
% -------------------------------------------------------------------

\subsection{A short excursus: why this class of models is non-identifiable}
% ~1/3 page. Side story that connects to the sigma_input reveal.

\subsubsection{The structural argument}
First-order kinetics (and richer topologies) let parameters interact so the
data constrain net fluxes out of a pool, not the individual splits.

\paragraph{Sketch}
For a linear pool system the observable stocks depend on the parameters only
through lumped combinations; distinct parameter sets give identical
predictions (parameter-space ridges). Two equations of placeholder math.
Full argument in Appendix~\ref{app:nonident}. [placeholder]

\paragraph{Signature in practice}
% [Bayes] — the lens: the Bayesian frame makes non-identifiability visible and
% honest (posterior vs prior, KL) rather than hiding it in a point estimate.
Prior-pinning, near-perfect anti-correlation ridges between same-pool
fractions, low information gain (KL of posterior vs prior). The Bayesian frame
is what lets us see this rather than hide it in a fitted point. Tie to Results.

\paragraph{Why it is not a defect to be ``fixed''}
It is intrinsic to the model class and the data content (few observations per
plot); it bounds what calibration can ever claim.

\paragraph{Managing it by convention, transparently}
We do not claim to resolve the rate$\leftrightarrow$fraction equifinality --- we
MANAGE it: where authoritative rate values exist (Yasso) we pin the chemical axis
and calibrate the partitioning; where they do not (simple models) we calibrate and
let the equifinality show. An explicit convention, not an identification claim ---
and the candour is what keeps this section credible rather than defensive. Ties
back to ``degrees of freedom set by external information'' in Methods.

\subsection{Structure is a weak lever in-sample, but it dominates the forecast}
% CAPSTONE of Side Arc A (its 2nd face) + bridge to the policy landing (the
% ensemble-spread map). Locks 2026-07-14/07-15: fold don't branch; ONE beat;
% plausibility = a CAUTIOUS lean to saturation (keep the Yasso humus-parking
% counter-caveat); policy angle = brief mention. Figure: F4(b) + the ensemble-SD map.
Within the calibrated window the six structures are near-interchangeable: complexity
buys no skill (Results), and the excursus just showed why --- the surplus degrees of
freedom are non-identified, with the inputs absorbing the misfit. Extrapolated, the
picture inverts. Run forward to 2084 the models agree where data constrain them and then
fan out sharply --- the single-pool SP1 climbs without bound ($\sim$133~tC\,ha$^{-1}$),
while the stiffer Yasso structures plateau ($\sim$108). The very structural choices the
data could not pin down now govern the projection. In-sample skill is blind to this: it
is scored only on the current stocks, which every model reproduces about equally well.

\subsubsection{A cautious read, not a ranking}
Process plausibility --- not fit --- gives a soft steer. Soils saturate, so SP1's
unbounded single-pool rise is the least physically credible trajectory, and the
saturating structures sit better with what we expect of a maturing sink. But the steer
is soft and we keep its counter-caveat explicit: under the fair frame Yasso carries the
most free degrees of freedom, so its plateau could be input-bound humus-parking rather
than a better-resolved dynamic. Present data cannot adjudicate between these readings.

\subsubsection{The honest output is the spread}
Because the divergence is real, decision-relevant, and unfalsifiable from present data,
the defensible product is not a chosen model but the \emph{ensemble spread} --- reported,
not hidden. This is the deepest reason the national maps carry a between-model
uncertainty panel beside the mean: the spatial SD is the forecast disagreement made
geographic. And the projection a greenhouse-gas inventory most needs --- multi-decade SOC
change --- is precisely the quantity in-sample skill cannot validate, which is why the
honest hand-off to policy is the spread, not a single number.
% Forward-pointer (future work on TESTING the divergence). Kept to 2 sentences.
Narrowing that spread will not come from more SOC campaigns --- two or three stocks per
plot cannot arbitrate multi-decade dynamics --- but from evidence that constrains the
\emph{trajectory}: long-term chronosequences and re-measured permanent plots, soil-warming
and litter-manipulation experiments, and process-based limits on how much carbon a maturing
boreal soil can hold. Until such constraints exist, the divergence is a genuine, reportable
uncertainty, not a modelling failure.

\subsection{A level offset that turned out to be ours, not the soil's}
% Persona/candor beat.
%
% ⚠ REWRITTEN 2026-09-02. The previous version asserted that a "slight, systematic
% over-prediction follows from the multiplicative-normal error model". The premise
% has since inverted twice over: (i) that artefact, taken alone, was shown to
% UNDER-predict, so the sign was wrong; and (ii) the uniform over-prediction it
% sought to explain NO LONGER EXISTS -- it closed when the likelihood moved to
% log-normal and the SOC target was rebuilt on a single homogenised basis. Kept as
% a candour beat because the episode is instructive, not because the bias survives.
An earlier stage of this work carried a uniform over-prediction: all six models
sat $17$--$20\%$ above the observations, by almost exactly the same margin,
despite being structurally unrelated. Uniformity across six independent
structures is not a biogeochemical signal; it points at something the six share.
They shared two things --- the fitting criterion and the target --- and both
turned out to be implicated. Replacing the multiplicative-normal likelihood with
a log-normal one removed the mechanism by which a prediction could widen its own
error bar, and rebuilding the SOC target on a single homogenised basis (coarse
fragments, layer definitions and campaign dating treated identically across
$1985$, $2006$ and $2024$) removed a stock inflation of comparable size.

What is left is small and no longer one-signed: median bias runs
$+1.0$ to $+3.3$~tC\,ha$^{-1}$ on the calibration plots and $-0.9$ to
$+0.7$~tC\,ha$^{-1}$ on the holdout, against stocks near $65$~tC\,ha$^{-1}$.
There is no longer a level offset to explain.

We report this because it is the kind of finding a model intercomparison is
built to produce and easy to misread. Had we not homogenised the target and the
likelihood first, the offset would have been available for attribution to soil
processes --- to a missing stabilisation mechanism, or to inputs --- in six
models at once, and the agreement between them would have looked like
corroboration. It was an artefact of the frame the six were being compared in.
The general point survives the specific one: in an intercomparison, a residual
common to \emph{every} structure is evidence about the comparison, not about the
soil.

\subsection{Limitations}

\subsubsection{The coarse pre-observation phase}
Linear litter interpolation under constant climate; high, deliberate
uncertainty.

\subsubsection{The flat-vs-rising litter tension}
Plot-level inputs vs.\ national reconstructions; present data cannot settle
it. State openly.

% -------------------------------------------------------------------
% [TEMP notes from the new sources (2026-07-16) — fold into prose later.]
\subsubsection{[TEMP] Between-campaign stock consistency (VMI8 vs Biosoil)}
\begin{itemize}
\item The apparent VMI8$\to$Biosoil sink ($\sim$184 gC\,m$^{-2}$\,yr$^{-1}$) is
  implausibly large; the paired 1985--2006 plots show the \emph{organic} layer
  stable but \emph{mineral} soil C nearly doubling (incl.\ deep 20--40\,cm) $\Rightarrow$
  a between-campaign mineral-soil measurement inconsistency, not accumulation.
\item Independent support: Peltoniemi et al.\ (2004) found \emph{mineral soil C does
  not change with stand age} (only the organic layer does), and a careful total of
  $6.8$\,kgC\,m$^{-2}$ to 1\,m ($\approx$ our VMI8 $\sim$63\,tC\,ha$^{-1}$; makes
  Biosoil/Komeetta $\sim$100+ look high). Cite for this caveat.
\item Handling (Methods): the 1985 stocks are \emph{down-weighted} via inflated
  observation variance ($\sim$2$\times$), not dropped; authoritative QC is ongoing
  (LUKE). Consequence: the ``$+26$ tC\,ha$^{-1}$ 1985 over-prediction'' is likely a
  low VMI8 measurement, not a model bias.
\end{itemize}
% -------------------------------------------------------------------

\subsection{A managed landscape out of equilibrium: Finnish forest history}
% [History] — the detailed examination promised in the Intro; justifies the
% non-equilibrium and the rising-input trajectory; feeds the reconstruction call.
The concrete history behind the violated equilibrium assumption: why Finnish
forest soils were accumulating, not at steady state, across the observation
period. Pays off the Intro's general statement.

\subsubsection{From exploitation to managed growth}
A sketch of the management eras --- heavy historical use, then regeneration and
intensified management, rising growing stock through the 20th century.
[placeholder --- sources to add]

\subsubsection{The imprint on litter inputs and stocks}
Rising biomass and shifting management make the input history non-stationary
--- precisely the signal an equilibrium start erases. Connect to the
$\sigma_{\mathrm{input}}$ reveal and the flat-vs-rising litter tension.

\subsubsection{Why this is the rule, not a Finnish peculiarity}
Most managed and recovering landscapes carry comparable histories; the lesson
generalizes. Tie back to the Intro's general statement.

\subsection{Outlook: toward a historical perspective on inputs}
The forward call. The real next step is reconstructing historical (and
belowground) input regimes.
% TWO-PRONGED FRONTIER, one lesson (add structure to the FORCING, not the kinetics):
%   (temporal) reconstruct the input HISTORY --- replace the linear 1917->1985 guess;
%   (spatial/class) stratify the input LEVEL by site class --- the red optional strand.
% The physical envelope UNIFIES both: it bounds sigma_init (the 1917/history input)
% AND sigma_input (the contemporary input) under one constraint, so ``reconstruct the
% input history'' literally means doing better than our linear guess at BOTH ends of
% the pre-run. Note: Komeetta 2024 = a 3rd SOC obs/plot -> won't rescue kinetic
% identifiability (12 fractions vs 3 points), but strengthens the input/init signal
% that the stratified strand leans on.

\subsubsection{From diagnosis to method (now folded into M\&M)}
% The physical input bound was the ``refinement'' beat --- it is now IMPLEMENTED as
% a Methods choice (M&M ``Physically-bounded input priors'' + App.~\ref{app:input-bounds}),
% so it is no longer future work. What remains genuinely FUTURE is the input HISTORY
% (temporal reconstruction) and belowground inputs.
The bound itself has moved to Methods; the remaining frontier is the input
\emph{history} (temporal) and belowground inputs. See Appendix~\ref{app:input-bounds}.

\subsubsection{An interdisciplinary, source-diverse effort}
Reconstruction needs evidence beyond the usual datasets --- grey literature
and non-traditional historical sources. Frame as a research program.

\subsubsection{Initialization as a calibratable, evidence-based step}
% [Bayes] — the thread closes: reconstructed history becomes the prior.
Replace the equilibrium default with priors grounded in reconstructed history,
calibrated under uncertainty. The thread closes here: better history $\to$
sharper prior $\to$ a start that is assumed less and inferred more. The
generalizable recommendation.

\paragraph{Concrete near-term step (to explore next session).}
% Use the NFI growing-stock history (F10b / storyline note Fig.1) as the pre-1985
% litter interpolation SHAPE, replacing the linear ramp. Litter scales with standing
% biomass, so the growing-stock curve (depleted mid-century -> rise) is a physically
% grounded, essentially free reconstruction of the input history -- turns the [History]
% thread into an actual model input and targets the coarse pre-1985 phase behind the
% +26 tC/ha 1985 over-prediction. The cheapest realisation of "reconstruct the input
% history"; a candidate improvement to the transient-init pre-run.
The transient-init pre-run currently ramps litter \emph{linearly} from the 1917
anchor to 1985. A near-term improvement is to use the \emph{shape} of the NFI
growing-stock history (\S[History]) as that interpolation function instead: since
litter scales with standing biomass, the growing-stock recovery curve is a
physically grounded, essentially cost-free reconstruction of the pre-observation
input path, and it targets the coarse pre-1985 phase directly.

% =====================================================================
%  OPTIONAL STRATIFIED-CALIBRATION STRAND  (red; drop entirely with \stratfalse)
%  Self-contained: removing it leaves the storyline whole, just smaller.
% =====================================================================
\ifstrat
{\color{red}
\subsection{[Optional strand] Locally stratified calibration (Yasso only)}
A droppable extension --- include only if data + time allow; the storyline stands
without it. It is the SPATIAL counterpart of the historical-reconstruction prong:
both add structure to the \emph{forcing}, never to the kinetics.

\subsubsection{Scope --- the operational model only}
Stratify Yasso (the operational model); SP1/TP2/TP3 stay global robustness anchors.
Consistent with the rate convention: Yasso already fixes its published rates, so its
free knobs (inputs + fractions) are the ones available to stratify. Triggered only
``if complexity emerges as needed'', not by default.

\subsubsection{What to stratify --- inputs and the humus-formation fraction}
The physically-bounded input $\sigma$ and the humus-formation fraction (the
stabilisation control) --- NOT the rates (fixed) nor the inter-pool splits
(non-identified). Both carry external, class-varying referents; humus formation is
a near-net stabilisation flux the total stock genuinely constrains, so of the
fractions it is the one worth stratifying.

\subsubsection{RF-guided selection, holdout-evaluated}
Use the residual random forest to CHOOSE the stratifying classes (its OOB $R^2$ is
the ceiling of covariate-recoverable local structure); then measure the stratified
model's gain on the HELD-OUT plots --- non-negotiable, or selecting and scoring on
the same data makes the improvement circular. Report ``RF says up to $X\%$ is
recoverable; the stratified mechanistic model recovers $Y\%$ on holdout.''

\subsubsection{Hierarchical, predictive-robustness framing}
Partial pooling: data-poor strata shrink to the global value (no minimum-$N$
problem); the between-stratum variance is itself a result. Headline = improved,
more ROBUST posterior predictions even where stratum parameters stay weakly
identified --- identifiability and predictive skill are distinct, and robustness is
what the GHG-inventory use needs.

\subsubsection{Equifinality as a measured quantity}
Bounded inputs vs.\ free humus formation compete for the local variance; the
physical bound PARTIALLY breaks their equifinality (one competitor has walls, the
other does not), so the learned variances softly attribute local variance to
input vs.\ stabilisation origin --- cross-checked by the RF driver identity
(productivity proxies $\to$ input; soil-chemistry proxies $\to$ stabilisation).
Frame as ``consistent with'', not proof. Expect modest holdout gains against a low
predictability ceiling --- modest-but-robust-and-honest beats an overclaimed jump.
}
\fi
% =====================================================================
%  end optional strand
% =====================================================================

\section{Conclusions}
% Written LAST; author shapes. Closes on: success -> Bayesian enabling ->
% refine with historical data collection.

\subsection{The experiment succeeded}
Transient, calibrated initialization let a model ladder reproduce the observed
accumulation and its incipient saturation --- the dynamic an equilibrium start
forecloses. State the proof of concept plainly.

\subsection{The Bayesian approach made it possible}
% [Bayes]
Treating the historical start as a calibrated parameter, with uncertainty
propagated into state and dynamics, is what turned the idea into a working
implementation. Credit the machinery without overclaiming novelty.

\subsection{Refinement needs historical data}
% [History]
The approach can go much further once the pre-observation history is
reconstructed from data rather than left to calibration --- the open frontier,
and the natural successor study. Close here.

% =====================================================================
\part{Appendices}
% Modular methods inserts. Each \subfile below is a standalone `subfiles`
% document under appendices/ that also compiles on its own. Comment a line to
% drop that module from the combined build; the full sigma_input module is
% written, the rest are structural stubs to be filled.
% =====================================================================
\appendix

\subfile{appendices/appendix_sigma_input}          % FULL — physically-bounded input priors
\subfile{appendices/appendix_prior_homogenization} % stub
\subfile{appendices/appendix_transient_init}       % stub
\subfile{appendices/appendix_nonidentifiability}   % stub
\subfile{appendices/appendix_exact_integration}    % stub
\subfile{appendices/appendix_delta_reconciliation} % FULL — SOC-change diagnostics reconciled
\subfile{appendices/appendix_rmse_posterior}       % FULL — per-draw RMSE, cross-run comparison

% ---------------------------------------------------------------------
\section*{Notes for collapsing (remove before submission)}
\begin{itemize}
  \item \textbf{Parts} = OCAR acts; will dissolve into standard
        Intro/Methods/Results/Discussion at cleanup.
  \item Many heading levels are intentional scaffolding; merge aggressively
        once the story is fixed.
  \item Keep the protagonist (initialization problem) constant; do not let
        $\sigma_{\mathrm{input}}$ or non-identifiability take over the spine.
\end{itemize}

\end{document}
